% !TEX root = ../main.tex
% ============================================================================
% Gluing Chapter, Part II:
%   Chart-wise gluing on toric CY$_3$ (Stratum i)
%
% First-principles elementary-derivation draft.  Every step is pinned to an
% explicit chart computation: Jacobian of the transition, action on the
% $(3,0)$-form, induced action on polyvectors, Fourier--Mukai multiplier,
% Koszul sign.  This draft is standalone (CG voice, no "Wave N" language,
% no process narration, no retraction bookkeeping); it subsumes the
% conifold two-chart construction of
% \texttt{chapters/examples/toric\_cy3\_coha.tex} \S conifold-two-chart-cech,
% feeds the McKay/orbifold Part III, and states the local-to-global
% theorem that terminates Stratum~i.
% ============================================================================

\providecommand{\bY}{\mathbb{Y}}

\part{Chart-wise gluing on toric CY$_3$ (Stratum i)}

\medskip

Stratum~i is the toric regime: every top-dimensional cone of the fan
contributes a $\C^3$-chart carrying the Schiffmann--Vasserot positive
half $Y^+(\widehat{\fgl}_1)$, and the global CoHA is the
\v{C}ech--$E_3$-factorisation-algebra homotopy colimit of the chart
data along the face lattice of the fan.  The simplest non-trivial
instance is the resolved conifold
\(
  \bY = \mathrm{Tot}(\cO_{\bP^1}(-1) \oplus \cO_{\bP^1}(-1))
\)
with its two-chart atlas $\bY = U_+ \cup U_-$.  The multi-chart
generalisation indexes over a toric fan; the compact-$4$-cycle
obstruction separates zero-interior-point diagrams from diagrams
requiring McKay resolution; the local-to-global theorem closes the
toric programme at the zero-interior-point boundary.

% ============================================================================
\section{\v{C}ech gluing of two charts: the conifold paradigm}
\label{sec:gluing-p2-conifold}
% ============================================================================

The obstruction resolved by a two-chart atlas is the absence of a global
$\C^3$-parametrisation: the resolved conifold $\bY$ is a rank-$2$ vector
bundle over $\bP^1$, not an affine variety.  Any chart-wise
Schiffmann--Vasserot construction must name the transition datum that turns
two copies of $\C^3$ into $\bY$, then transport that datum through every
intermediate step --- the $(3,0)$-form, the polyvector fields, the Hall
product, the super structure.

\subsection{The two-chart atlas}

Let $z \in \bP^1$ be the standard affine coordinate on the base, with
$\tilde z = z^{-1}$ on the complementary chart.  On the total space
\(
  \bY = \mathrm{Tot}(\cO_{\bP^1}(-1) \oplus \cO_{\bP^1}(-1))
\)
introduce fibre coordinates $u_+, v_+$ in the trivialisation over the
$z$-chart and $u_-, v_-$ in the trivialisation over the $\tilde z$-chart.
The cover
\[
  U_+ \;=\; \{z \neq \infty\} \times \C^2 \;\cong\; \C^3,
  \qquad
  U_- \;=\; \{\tilde z \neq \infty\} \times \C^2 \;\cong\; \C^3,
\]
has overlap $U_+ \cap U_- = (\bP^1 \setminus \{0, \infty\}) \times \C^2
\cong \bC^\times \times \C^2$.  A section of $\cO_{\bP^1}(-1)$ is a function
$f(z)$ on $U_+$ that transforms to $\tilde z f(\tilde z^{-1}) = z^{-1} f(z)$
on $U_-$.  Writing $u_\pm, v_\pm$ as the fibre coordinates of the two $\cO(-1)$
summands one finds the transition
\begin{equation}\label{eq:conifold-transition-chart}
  \tilde z \;=\; z^{-1},
  \qquad
  u_- \;=\; z\, u_+,
  \qquad
  v_- \;=\; z\, v_+,
\end{equation}
as the natural datum forced by the $\cO(-1) \oplus \cO(-1)$ twist.  The
standard atlas of $\bP^1$ carries an orientation flip across
$z \leftrightarrow \tilde z$: the holomorphic $1$-form $dz = -\tilde z^{-2}\,
d\tilde z$ reverses sign through the change of chart, and this sign
propagates to the $(3,0)$-form computation of the next subsection.

\subsection{Jacobian of the transition on the $(3,0)$-form}

The Calabi--Yau $(3,0)$-form $\Omega_{\bY}$ is specified on $U_+$ by the
requirement that it trivialise the canonical bundle of $U_+ \cong \C^3$:
\[
  \Omega_{\bY}\big|_{U_+} \;=\; dz \wedge du_+ \wedge dv_+.
\]
To transport $\Omega_{\bY}$ to $U_-$ one computes the Jacobian of the
transition~\eqref{eq:conifold-transition-chart} step by step.  The exterior
derivatives are
\[
  dz \;=\; -\tilde z^{-2}\, d\tilde z,
  \qquad
  du_+ \;=\; d(z^{-1} u_-)
          \;=\; -z^{-2}\, u_-\, dz + z^{-1}\, du_-,
  \qquad
  dv_+ \;=\; -z^{-2}\, v_-\, dz + z^{-1}\, dv_-,
\]
so that
\[
  dz \wedge du_+ \;=\; dz \wedge z^{-1}\, du_-,
  \qquad
  dz \wedge du_+ \wedge dv_+ \;=\; dz \wedge z^{-1}\, du_- \wedge z^{-1}\, dv_-
  \;=\; z^{-2}\, dz \wedge du_- \wedge dv_-.
\]
Substituting $dz = -\tilde z^{-2}\, d\tilde z$ and $z^{-2} = \tilde z^{\,2}$:
\begin{equation}\label{eq:jacobian-calculation}
  \Omega_{\bY}\big|_{U_+}
  \;=\; dz \wedge du_+ \wedge dv_+
  \;=\; z^{-2}\, dz \wedge du_- \wedge dv_-
  \;=\; (\tilde z^{\,2}) \cdot (-\tilde z^{-2})\, d\tilde z \wedge du_- \wedge dv_-
  \;=\; -d\tilde z \wedge du_- \wedge dv_-.
\end{equation}
The fibre Jacobian $z^{-2} = \tilde z^{\,2}$ cancels the cotangent factor
$-\tilde z^{-2}$ from $dz$ exactly; the surviving sign is the
base-orientation flip of $\bP^1$.  This is the statement that $K_{\bY}$
is trivial as a line bundle on $\bY$ --- its transition cocycle on the
overlap $U_+ \cap U_-$ collapses to $\pm 1$ --- with the base-orientation
sign $-1$ read off the global trivialisation~\eqref{eq:jacobian-calculation}.
The fibre twist $z^{-2}$ (two factors of $\cO(-1)$, one from each
$\cO_{\bP^1}(-1)$ summand) compensates the cotangent factor $z^{-2}$ on
the base $1$-form $dz$, leaving the constant $-1$ as the whole transition
multiplier on $\Omega_\bY$.

\subsection{The chart-wise Schiffmann--Vasserot CoHA}

Each chart $U_\pm \cong \C^3$ carries the Jordan-triple quiver with
potential
\[
  Q_{\mathrm{Jor}} : \;\; \underbrace{\bullet}_{0}
  \;\;\text{with three loops}\;\; x_\pm, y_\pm, z_\pm,
  \qquad
  W_\pm \;=\; \mathrm{tr}\bigl(x_\pm[y_\pm, z_\pm]\bigr).
\]
The critical CoHA of $(Q_{\mathrm{Jor}}, W_\pm)$ is the positive half of the
affine Yangian of $\widehat{\fgl}_1$ (Schiffmann--Vasserot 2012
\texttt{arXiv:1202.2756}; Tsymbaliuk 2014 \texttt{arXiv:1404.5240}),
\[
  \CoHA\bigl(U_\pm, W_\pm\bigr) \;\cong\; Y^+\bigl(\widehat{\fgl}_1\bigr),
\]
with generators $e_r^{(\pm)}, f_r^{(\pm)}, h_r^{(\pm)}$ indexed by $r \in
\Z_{\geq 0}$ and structure functions given by the RSYZ shuffle kernel
(Rapcak--Soibelman--Yang--Zhao 2020 \texttt{arXiv:1810.10402} \S5,
equation (5.3), specialised to the single-vertex Jordan-triple case).
At weight $(1,1)$ the shuffle kernel evaluates to
\begin{equation}\label{eq:structure-constant-1-1}
  \varphi_{11}(u_1; u_2) \;=\; \frac{1}{(u_1 - u_2 - \varepsilon_1)(u_1 - u_2 - \varepsilon_2)}
  \;\;\xrightarrow{\;u_2 \to u_1\;}\;\;
  \frac{1}{\varepsilon_1 \varepsilon_2},
\end{equation}
the structure constant that every downstream computation must reproduce.

\subsection{Fourier--Mukai transition kernel}

On the overlap $U_+ \cap U_- = \bC^\times \times \C^2$ the derived category
$D^b(\mathrm{Coh}(U_+ \cap U_-))$ carries the autoequivalence implementing
the transition~\eqref{eq:conifold-transition-chart} at the level of sheaves.
Its integral kernel is
\begin{equation}\label{eq:FM-kernel-conifold}
  K_{+-} \;=\; \cO_\Delta \otimes \pi_1^*\,\cO_{\bP^1}(-1)
  \;\in\; D^b\bigl(\mathrm{Coh}\,(U_+ \cap U_-) \times (U_+ \cap U_-)\bigr),
\end{equation}
with $\Delta$ the diagonal and $\pi_1$ the projection onto the first factor.
The transform $T_{+-}(\cF) := \pi_{2,*}\bigl(\pi_1^*\cF \otimes K_{+-}\bigr)$
acts on polyvector fields
$\mathrm{PV}^\bullet(U_\pm) = \bigoplus_p \bigwedge^p T_{U_\pm}$ through the
chain rule applied to~\eqref{eq:conifold-transition-chart} combined with the
$\cO_{\bP^1}(-1)$-twist.

\paragraph{Chain rule on the cotangent basis.}
Differentiate~\eqref{eq:conifold-transition-chart} to obtain
\[
  du_+ \;=\; d(z^{-1} u_-) \;=\; -z^{-2}\, u_-\, dz + z^{-1}\, du_-,
  \qquad
  dv_+ \;=\; -z^{-2}\, v_-\, dz + z^{-1}\, dv_-.
\]
On the vertical (fibre-degree) sector --- the sector entering the
Schiffmann--Vasserot Hall product --- only the $z^{-1}\,du_-$ and
$z^{-1}\,dv_-$ terms survive (the $dz$-terms drop under pairing with the
Jordan-triple fibre polyvectors by~\eqref{eq:structure-constant-1-1}).
The cotangent multiplier is therefore $z^{-1}$ per fibre slot, giving
$z^{-(a+b)}$ on a fibre-form $du_+^a \wedge dv_+^b$.

\paragraph{Polyvector transport via $\Omega_\bY$-duality.}
The critical CoHA identifies polyvectors with differential forms via
Serre duality against the $(3,0)$-form $\Omega_\bY$:
\[
  \Lambda^p\, T_{\bY} \;\cong\; \Omega^{3-p}_\bY \otimes K_\bY^{-1}
  \;\cong\; \Omega^{3-p}_\bY
\]
using triviality of $K_\bY$ (equation~\eqref{eq:jacobian-calculation}).
Under this identification the fibre-polyvector $\partial_{u_+}^a
\partial_{v_+}^b$ corresponds (up to sign) to the fibre-form
$du_+^a \wedge dv_+^b$ contracted against the base $dz$-direction of
$\Omega_\bY$.  Its transport under $T_{+-}$ is the cotangent transformation
$du_+^a \wedge dv_+^b = z^{-(a+b)}\, du_-^a \wedge dv_-^b$.

\paragraph{$\cO(-1)$-kernel twist.}
The factor $\pi_1^*\cO_{\bP^1}(-1)$ in $K_{+-}$ contributes a further
multiplier $z^{-1}$ on the scalar sector: a section of
$\cO_{\bP^1}(-1)$ on $U_+$ written as $g(z)$ is expressed on $U_-$ as
$\tilde z\, g(\tilde z^{-1}) = z^{-1}\, g(z)$.

\paragraph{Net multiplier.}
Combining the cotangent re-section and the kernel twist,
\[
  \bigl[z^{-(a+b)}\bigr]_{\text{cotangent chain rule}}
    \;\cdot\;
  \bigl[z^{-1}\bigr]_{\cO(-1)\text{-kernel}}
  \;=\;
  z^{-(1 + a + b)}.
\]
The Fourier--Mukai transport of a Jordan-triple fibre polyvector therefore
reads
\begin{equation}\label{eq:FM-multiplier}
  T_{+-}^* :\;
  f(z, u_+, v_+)\, \partial_{u_+}^a\, \partial_{v_+}^b
  \;\longmapsto\;
  f(\tilde z^{-1}, \tilde z\, u_-, \tilde z\, v_-)\;
  z^{-(1 + a + b)}\;
  \partial_{u_-}^a\, \partial_{v_-}^b.
\end{equation}
The base-polyvector sector (powers of $\partial_z$) picks up the
base-orientation multiplier $(-\tilde z^{\,2})$ per $\partial_z$-slot; the
Jordan-triple selection rule~\eqref{eq:structure-constant-1-1} restricts
the Hall product to the fibre sector $\alpha = 0$ on which the base
orientation enters only through the global sign
of~\eqref{eq:jacobian-calculation}.

\subsection{Homotopy colimit and super-Yangian output}

The pair $(U_+, U_-)$ with overlap $U_+ \cap U_-$ forms a two-object diagram
in the category of $\C^3$-charts.  The chart-wise factorisation algebras
$\cF_\pm = \Phi^{FA}_3(\Gamma_\pm)$ (Ginzburg dgas of the Jordan-triple
quiver with potential, \S\ref{sec:gluing-p2-conifold}; Kontsevich--Tamarkin
$E_3$-formality + Costello--Li holomorphic locality) glue by the homotopy
colimit
\begin{equation}\label{eq:homotopy-colimit}
  \cF_{\bY} \;=\; \mathrm{hocolim}\Bigl(
    \cF_+\big|_{U_+ \cap U_-}
    \;\xrightarrow{\;T_{+-}\;}\;
    \cF_-\big|_{U_+ \cap U_-}
  \Bigr),
\end{equation}
an $E_3$-factorisation algebra on $\bY$ whose cohomology is
$H^\bullet(\cF_{\bY}) = \mathrm{PV}^\bullet(\bY)$.  Under the CoHA assignment
the homotopy colimit produces a new associative algebra whose generators
split into two families indexed by the two charts and whose structure
constants are determined by the multiplier~\eqref{eq:FM-multiplier}.

\paragraph{The super structure as a Koszul sign.}
The base-orientation flip in $\Omega_{\bY}|_{U_-} = -d\tilde z \wedge du_-
\wedge dv_-$ (equation~\eqref{eq:jacobian-calculation}) contributes a
Koszul sign $\sigma = -1$ on cross-chart Hall products: the orientation of
the measure against which the critical CoHA integrates polyvectors
reverses between charts, so bilinear pairings $\langle \cdot, \cdot
\rangle_\pm$ that glue across the overlap pick up a factor of $-1$.  In
the chart-wise shuffle presentation this is the parity assignment of
the Schiffmann--Vasserot generators across the two charts, producing
super-Yangian-type anticommutators on the combinations
\[
  E^{(0)} \;=\; e_0^{(+)} + e_0^{(-)},
  \qquad
  E^{(1)} \;=\; e_0^{(+)} - e_0^{(-)},
\]
where the even combination $E^{(0)}$ generates a bosonic subalgebra and the
odd combination $E^{(1)}$ generates a fermionic sector with
\[
  \bigl\{E^{(1)}, E^{(0)}\bigr\}_+ \;=\; K^{(0)},
\]
the central element extracted from the residue at $z_1 = z_2$ of the
chart-wise Jordan shuffle difference~\eqref{eq:structure-constant-1-1}.
The two-chart assembly therefore yields
\begin{equation}\label{eq:conifold-master}
  \CoHA(\bY) \;\cong\; Y^+\bigl(\widehat{\fgl}(1|1)\bigr)^{\mathrm{con}},
\end{equation}
the positive half of the conifold affine super-Yangian.  The super
structure on $\widehat{\fgl}(1|1)$ arises from the Koszul sign
$\sigma = -1$ of the base-orientation flip~\eqref{eq:jacobian-calculation}
combined with the isotropy $(\alpha_0, \alpha_0) = (\alpha_1, \alpha_1) = 0$
(Kac 1977 Thm.~2.4 on odd simple roots).

\paragraph{Three-route verification of the $(1,1)$ structure constant.}
The $(1,1)$ structure constant of $\CoHA(\bY)$ equals $(\varepsilon_1
\varepsilon_2)^{-1}$ through three independent routes:
\begin{enumerate}[label=\emph{(\roman*)}, leftmargin=2em]
\item \emph{\v{C}ech two-chart gluing}: the limit $u_2 \to u_1$ of the
  chart-wise Jordan shuffle~\eqref{eq:structure-constant-1-1}.
\item \emph{Van den Bergh tilting}: the integral of the $\varphi_W$-twisted
  Euler class over $\mathrm{Ext}^1(S_0, S_1) = \C^2$, the rank-$2$
  Ext-space between the simple modules of the Klebanov--Witten quiver.
\item \emph{Klebanov--Witten super-shuffle}: the Jordan kernel residue at
  $z_1 = z_2$ of the direct super-shuffle algebra.
\end{enumerate}
Agreement across the three routes pins $(\varepsilon_1 \varepsilon_2)^{-1}$
as a lane-independent invariant, confirming~\eqref{eq:conifold-master} at
the coefficient level.

% ============================================================================
\section{Multi-chart toric fan generalisation}
\label{sec:gluing-p2-multi-chart}
% ============================================================================

The two-chart assembly generalises to any local toric CY$_3$.  The obstruction
resolved is the absence of a global affine chart: a toric variety is covered
by affine opens $U_\sigma$ indexed by the maximal cones $\sigma$ of the fan,
and the global CoHA is assembled by running the \v{C}ech machinery along the
face lattice of the fan.

\subsection{Toric fan and chart decomposition}

Let $X = X_\Sigma$ be a smooth toric CY$_3$ with fan $\Sigma$ in the lattice
$N \cong \Z^3$.  The Calabi--Yau condition is that the primitive generators
of all rays of $\Sigma$ lie in a common affine hyperplane, equivalently
that there exists a lattice vector $m \in M = N^\vee$ such that
$\langle m, v_\rho \rangle = 1$ for every ray generator $v_\rho$.  A
natural choice places all ray generators at height $1$ along the
last coordinate; the toric diagram is then the $2$-dimensional polytope
obtained by projecting the rays to the plane $\{x_3 = 1\}$.

\paragraph{Chart structure at top cones.}
Each maximal (top-dimensional) cone $\sigma \in \Sigma(3)$ of a smooth
toric CY$_3$ is generated by exactly three rays
$\rho_1, \rho_2, \rho_3$ forming a $\Z$-basis of $N$, and determines an
affine chart
\[
  U_\sigma \;=\; \mathrm{Spec}\,\C[\sigma^\vee \cap M] \;\cong\; \C^3.
\]
Coordinates $(x_1^\sigma, x_2^\sigma, x_3^\sigma)$ on $U_\sigma$ are the
characters of the dual basis of $\sigma^\vee$.  The Jordan-triple
potential on each chart is
\[
  W_\sigma \;=\; \mathrm{tr}\bigl(x_1^\sigma [x_2^\sigma, x_3^\sigma]\bigr),
\]
and the chart-wise CoHA is $\CoHA(U_\sigma, W_\sigma) \cong
Y^+(\widehat{\fgl}_1)$ by the Schiffmann--Vasserot theorem applied to each
$\C^3$ with the standard Jordan-triple.

\paragraph{Equivariant weights and the CY condition.}
The three-torus $T^3 = (\bC^\times)^3$ acts on each $U_\sigma \cong \C^3$
with equivariant weights $(\varepsilon_1^\sigma, \varepsilon_2^\sigma,
\varepsilon_3^\sigma)$.  The Calabi--Yau constraint on $W_\sigma$ to be a
tri-vector of vanishing weight is
\[
  \varepsilon_1^\sigma + \varepsilon_2^\sigma + \varepsilon_3^\sigma \;=\; 0,
\]
which specialises the chart-wise weights to those of a local CY$_3$.  Along
a compact direction (a ray $\rho$ in $\Sigma(1)$ whose dual edge in the
toric diagram is an internal edge) the equivariant parameter vanishes:
the two fixed points of the compact $\bP^1$-base on that edge would carry
opposite weights, obstructing global regularity of sections.  This is the
mechanism that set $\varepsilon_3 = 0$ for the conifold base in
\S\ref{sec:gluing-p2-conifold}.

\subsection{Gluing cocycle on the face lattice}

Two top cones $\sigma, \sigma'$ of $\Sigma(3)$ sharing a common face
$\tau = \sigma \cap \sigma' \in \Sigma(2)$ of codimension one define an
overlap $U_\sigma \cap U_{\sigma'} = U_\tau$.  The face $\tau$ is a
two-dimensional cone, generated by two rays; the corresponding overlap is
\[
  U_\tau \;\cong\; \bC^\times \times \C^2,
\]
with the $\bC^\times$-factor encoding the direction transverse to $\tau$
within the full three-dimensional lattice.  The transition between the chart coordinates
$(x_1^\sigma, x_2^\sigma, x_3^\sigma)$ and $(x_1^{\sigma'}, x_2^{\sigma'},
x_3^{\sigma'})$ follows the lattice change of basis between $\sigma^\vee$
and $(\sigma')^\vee$: when $\sigma'$ differs from $\sigma$ by a toric flop
sending ray $\rho_1$ to $\rho_1'$ across the shared face $\tau$ (generated
by $\rho_2, \rho_3$), the conifold-type transition is
\begin{equation}\label{eq:multi-chart-transition}
  x_1^{\sigma'} \;=\; (x_1^\sigma)^{-1},
  \qquad
  x_i^{\sigma'} \;=\; x_1^\sigma \cdot x_i^\sigma
  \;\; \text{for } i = 2, 3,
\end{equation}
in direct analogy with the conifold
transition~\eqref{eq:conifold-transition-chart}.  More generally, if the
toric flop twists the fibre $\rho_1$-direction with multiplicity $(k_2, k_3)$
along $\rho_2, \rho_3$ (equivalently, the fibre bundle over the shared
$\bP^1_\tau$ is $\cO(-k_2) \oplus \cO(-k_3)$), the transition reads
$x_i^{\sigma'} = (x_1^\sigma)^{k_i}\, x_i^\sigma$ on the $i$-th fibre
coordinate.  The Fourier--Mukai kernel on the overlap is
\[
  K_{\sigma \sigma'} \;=\; \cO_\Delta \otimes \pi_1^* \cO_{\bP^1_{\tau}}(-1),
\]
the $\cO(-1)$-twist sourced by the compact $\bP^1_\tau$-base transverse
to $\tau$; the polyvector multiplier
$(x_1^\sigma)^{-(1 + k_2\, a + k_3\, b)}$ generalises~\eqref{eq:FM-multiplier}
with $a, b$ the fibre polyvector indices.  The Jacobian
argument of \S\ref{sec:gluing-p2-conifold} carries over with $z$ replaced
by $x_1^\sigma$ on each codimension-one face, and with the fibre twist
exponents $(k_2, k_3)$ replacing the conifold's $(1, 1)$; the
Calabi--Yau condition $k_2 + k_3 = 2$ (sum of $\cO(-k_i)$-degrees equals
$-\deg K_{\bP^1_\tau} = -2$) enforces triviality of
$K_{X}|_{\bP^1_\tau}$.

\paragraph{The face lattice cocycle.}
The collection $\{K_{\sigma \sigma'}\}_{\sigma, \sigma' \text{ adjacent}}$
forms a cocycle on the face lattice of $\Sigma$: at any triple overlap
$U_\sigma \cap U_{\sigma'} \cap U_{\sigma''}$ (corresponding to a
codimension-two face $\tau = \sigma \cap \sigma' \cap \sigma''$), the
compatibility
\[
  K_{\sigma \sigma''} \;=\; K_{\sigma' \sigma''} \circ K_{\sigma \sigma'}
\]
holds up to canonical isomorphism, because each kernel is the pullback of
$\cO_{\bP^1}(-1)$ along a $\bP^1$-fibre direction and the fibre directions
compose additively in the toric fan.  The cocycle class is the derivative
of the data on the fan's $1$-skeleton and is classified by the sheaf
cohomology group $H^1(\Sigma, \underline{\mathrm{Pic}}(\bP^1))$ on the fan
viewed as a CW complex.

\subsection{Galakhov--Li--Yamazaki bond factors}

Galakhov--Li--Yamazaki (2021 \texttt{arXiv:2106.01230}; Li--Yamazaki 2020
\texttt{arXiv:2003.08909} eq~(8.125)) present the multi-chart gluing data
via \emph{bond factors} $\varphi^{a \Rightarrow b}(u)$ indexed by ordered
pairs $(a, b)$ of vertices in the toric quiver.  The convention is:
\begin{equation}\label{eq:gly-bond-factor}
  \varphi^{a \Rightarrow b}(u)
  \;=\;
  \frac{\prod_{I \in \mathrm{Arr}(a \to b)} (u + h^{a \to b}_I)}{%
        \prod_{J \in \mathrm{Rel}(a \to b)} (u + h^{a \to b}_J)},
\end{equation}
with $h^{a\to b}_I$ the equivariant weight of the $I$-th arrow from
vertex $a$ to vertex $b$, and $h^{a\to b}_J$ the equivariant weight of
the $J$-th Jacobi relation ($\partial W / \partial (\text{arrow})$) tying
arrows from $a$ into and out of a trip through $b$.  The \emph{numerator
is arrow weights; the denominator is Jacobi-relation weights}; the
$(u + h)$-convention signs track the Drinfeld shifted-$R$-matrix
normalisation (Neguț 2015 \texttt{arXiv:1505.01528} eq~(1.6);
Tsymbaliuk 2014 \texttt{arXiv:1404.5240}).

\paragraph{Conifold bond factor from first principles.}
The conifold quiver has two vertices $\{0, 1\}$, two arrows
$a_1, a_2 : 0 \to 1$ with equivariant weights $h_1, h_2$, two arrows
$b_1, b_2 : 1 \to 0$ with weights $-h_1, -h_2$, and the Klebanov--Witten
superpotential $W_{\mathrm{con}} = a_1 b_1 a_2 b_2 - a_1 b_2 a_2 b_1$.
The Jacobi relations $\partial W / \partial b_i = 0$ and $\partial W /
\partial a_j = 0$ pair arrows from $0 \to 1$ into and out of cycles
through vertex $1$; for the bond $0 \Rightarrow 1$, the two Jacobi
relations carry weights $0$ and $h_1 + h_2$ (derived by summing the
weights of the arrows meeting each relation-index).  Substituting into
the GLY formula~\eqref{eq:gly-bond-factor}:
\begin{equation}\label{eq:conifold-bond-factor}
  \varphi^{0 \Rightarrow 1}_{\mathrm{con}}(u)
  \;=\;
  \frac{(u + h_1)(u + h_2)}{u \cdot (u + h_1 + h_2)},
\end{equation}
matching Neguț's conifold shuffle kernel (Neguț 2015
\texttt{arXiv:1505.01528} eq~(1.6)) and Li--Yamazaki's
BPS-quiver-algebra master formula (\texttt{arXiv:2003.08909} eq~(8.125)).
At the Calabi--Yau torus-equivariant slice $\varepsilon_1 + \varepsilon_2
+ \varepsilon_3 = 0$, the conifold weights are $h_1 = \varepsilon_1$,
$h_2 = \varepsilon_2$, and $h_1 + h_2 = \varepsilon_1 + \varepsilon_2 =
-\varepsilon_3$; substituting into~\eqref{eq:conifold-bond-factor}:
\[
  \varphi^{0 \Rightarrow 1}_{\mathrm{con}}(u)
  \;=\;
  \frac{(u + \varepsilon_1)(u + \varepsilon_2)}{u\,(u + \varepsilon_1 + \varepsilon_2)}
  \;=\;
  \frac{(u + \varepsilon_1)(u + \varepsilon_2)}{u\,(u - \varepsilon_3)}.
\]
Both presentations are equivalent through the Calabi--Yau constraint
$\varepsilon_3 = -\varepsilon_1 - \varepsilon_2$.

\paragraph{Arrow-weight / relation-weight trace on the toric diagram.}
The arrow count $|\mathrm{Arr}(a \to b)|$ matches the number of toric
diagram edges joining the dual vertices of $a$ and $b$; the
Jacobi-relation count $|\mathrm{Rel}(a \to b)|$ matches the number of
toric triangles sharing the edge.  For the conifold ($2$ arrows, $2$
relations per direction) the $2 + 2$ tally reflects the two-edge,
two-triangle diagram.  For local $\bP^2$ the bond factors are labelled
by the three $\Z_3$-characters of the McKay quiver (Beilinson nine-arrow
structure, \S\ref{sec:gluing-p2-compact-cycle}); for the banana
threefold, by the three edges of the banana diagram; for the suspended
pinch point, by the pinch-point bond structure.

\subsection{Examples}

\paragraph{$\C^3/\Z_n$ (linear toric).}
The fan has $n$ maximal cones arrayed along the toric diagram edge, each
$\C^3$; the face lattice is a linear chain with $n - 1$ codimension-one
faces; the gluing cocycle is a chain of $n - 1$ FM kernels, each a
$\cO(-1)$-twist.  The multi-chart \v{C}ech assembly produces the
$A_{n-1}$-type McKay CoHA, matching the Beasley--Plesser antisymmetric
quiver on $\C^3/\Z_n$.

\paragraph{Banana threefold.}
Three ray clusters meeting at a common vertex produce a fan with three
maximal cones and three codimension-one faces.  The \v{C}ech assembly
requires three FM kernels on a triangular face lattice; the resulting
CoHA is a triangular gluing of three Schiffmann--Vasserot
$Y^+(\widehat{\fgl}_1)$ copies, governed by the three-arrow banana quiver
with cubic potential.

\paragraph{Suspended pinch point.}
Four maximal cones sharing a common pinch ray produce a fan with four
charts and six codimension-one faces.  The \v{C}ech cocycle is defined on
a square face lattice with one internal pinch; the CoHA assembly
reproduces the SPP quiver of Aganagic--Karch--L\"ust--Miemiec (1999).

% ============================================================================
\section{The compact $4$-cycle obstruction}
\label{sec:gluing-p2-compact-cycle}
% ============================================================================

The obstruction resolved here is subtle: a local toric CY$_3$ with a
compact divisor --- a maximal cone whose dual vertex lies in the
\emph{interior} of the toric diagram --- admits a chart-wise $\C^3$ cover,
but the direct \v{C}ech--$E_3$-factorisation assembly of
\S\ref{sec:gluing-p2-multi-chart} misses the cohomology class of the
compact divisor.

\subsection{Interior lattice points and compact divisors}

Let $X_\Sigma$ be a local toric CY$_3$ with toric diagram $\Delta_\Sigma$ a
lattice polygon in the plane $\{x_3 = 1\}$.  The duality between toric
fans and toric diagrams identifies:
\[
  \text{rays of } \Sigma \;\longleftrightarrow\; \text{vertices of } \Delta_\Sigma,
  \qquad
  \text{maximal cones} \;\longleftrightarrow\; \text{internal triangulation triangles},
\]
and one distinguishes:
\begin{itemize}
\item \emph{boundary rays}: rays whose dual vertex lies on the boundary of
  $\Delta_\Sigma$; these correspond to non-compact divisors in $X$;
\item \emph{interior rays}: rays whose dual vertex lies in the interior
  of $\Delta_\Sigma$; these correspond to compact divisors, i.e.\ compact
  $4$-cycles in $X$.
\end{itemize}
The resolved conifold has no interior lattice point: its toric diagram is a
unit square with vertices $(0,0), (1,0), (0,1), (1,1)$ and no interior
lattice point, consistent with $\chi(\bY) = 2$ and no compact $4$-cycle.
Local $\bP^2 = \mathrm{Tot}(K_{\bP^2})$ has exactly one interior lattice
point: its toric diagram is the triangle with vertices $(1,0), (0,1),
(-1,-1)$ (the three non-compact divisor rays at height $1$) and the
interior lattice point $(0,0)$ at the origin; the ray $(0,0,1)$ attached
to this interior point generates the zero-section $\bP^2 \subset K_{\bP^2}$,
a compact $4$-cycle.

\subsection{Why chart-wise gluing misses the compact divisor}

The chart-wise \v{C}ech assembly of \S\ref{sec:gluing-p2-multi-chart} has
three ingredients on each chart: the Schiffmann--Vasserot
$Y^+(\widehat{\fgl}_1)$, the equivariant torus action, and the FM kernel
on codimension-one overlaps.  None of these detect the cohomology of a
compact divisor:

\begin{itemize}
\item The chart-wise polyvector complex $\mathrm{PV}^\bullet(U_\sigma)$ is
  concentrated in $(0, *)$-degree (holomorphic polyvector fields on
  $\C^3$), and has no contribution from the class of a compact surface.
\item The FM kernels $K_{\sigma \sigma'}$ transport polyvectors across
  codimension-one faces but do not generate new cohomology classes:
  they are sheaves of derived Morita data, not cycle classes.
\item The $T^3$-equivariant localisation concentrates all integrals at
  the fixed points of the torus action, i.e.\ the maximal cones; the
  contribution of a compact divisor is the \emph{sum over the fixed
  points contained in the divisor}, which must be computed by a separate
  procedure.
\end{itemize}
The Neguț wheel-condition corrections --- the shuffle-algebra
relations forced by the presence of an interior lattice point --- are
missed by the direct \v{C}ech assembly.  For local $\bP^2$ with single
interior point at the origin of the triangulation, the correct shuffle
algebra is the $\Z_3$-McKay variant of $Y^+(\widehat{\fgl}_1)$ whose
kernel carries an additional wheel-condition multiplier tracking the
$\Z_3$-orbit of the interior lattice point (Neguț 2015
\texttt{arXiv:1512.06473}; Rapcak--Soibelman--Yang--Zhao 2020
\texttt{arXiv:2001.10549}).  The direct multi-chart \v{C}ech assembly of
\S\ref{sec:gluing-p2-multi-chart} picks up only the boundary contributions
to $\varphi_{a \Rightarrow b}$ and misses this wheel-condition
multiplier.

\subsection{Resolution via McKay/orbifold route}

The standard resolution of the obstruction is to pass to the orbifold
quotient: local $\bP^2 = [\C^3/\Z_3]^{\mathrm{res}}$ (the crepant
resolution of the McKay orbifold), and apply the Bridgeland--King--Reid
equivalence
\[
  D^b\bigl(\mathrm{Coh}(K_{\bP^2})\bigr) \;\simeq\; D^b\bigl(\mathrm{Coh}^{\Z_3}(\C^3)\bigr)
\]
(Bridgeland--King--Reid 2001 \texttt{arXiv:math/9908027} Thm.~1.2).  The
right-hand category is a single $\C^3$-chart with $\Z_3$-equivariance;
the CoHA of this orbifold chart carries the Neguț wheel-condition
multipliers automatically, since the $\Z_3$-action on $\C^3$ produces
exactly the interior-point shifts at each orbit element.  The compact
$4$-cycle cohomology is recovered as the $\Z_3$-fixed part of the
$\C^3/\Z_3$ inertia stack.

Alternatively, one may keep the multi-chart cover and enrich each
chart-wise shuffle kernel with a higher-rank Neguț correction: replace
$\varphi_{a \Rightarrow b}(u)$ by the shuffle kernel
$\varphi_{a \Rightarrow b}^{\mathrm{rank-}r}(u)$ of the rank-$r$ shuffle
algebra, with $r$ the number of interior lattice points of the toric
diagram.  The two resolutions --- McKay/orbifold route (Part III) and
higher-rank shuffle correction --- are Morita-equivalent under the
Rapcak--Soibelman--Yang--Zhao toric theorem (2020
\texttt{arXiv:2001.10549} Thm.~B), but differ in their chart-wise
presentations.

\begin{proposition}[Compact-$4$-cycle obstruction]
\label{prop:compact-cycle-obstruction}
Let $X_\Sigma$ be a local toric CY$_3$ with toric diagram $\Delta_\Sigma$
carrying $r \geq 1$ interior lattice points.  The direct multi-chart
\v{C}ech assembly of \S\ref{sec:gluing-p2-multi-chart} reproduces
$\CoHA_{\mathrm{DM}}(X_\Sigma)$ only after replacing the chart-wise
single-vertex Jordan-triple data by either:
\begin{enumerate}[label=\emph{(\alph*)}, leftmargin=2em]
\item the $G$-equivariant Schiffmann--Vasserot data on $[\C^3 / G]$,
  where $G \subset \mathrm{SU}(3)$ is the McKay group whose crepant
  resolution $\widetilde{\C^3/G}$ recovers $X_\Sigma$ (McKay/orbifold
  route; applies when $\Delta_\Sigma$ is a McKay triangle, including
  $X_\Sigma = K_{\bP^2}$ with $G = \Z_3$);
\item the multi-vertex Neguț shuffle algebra $Y^+(Q_X, W_X)$ whose
  vertex set matches the full toric quiver of $X_\Sigma$ (interior and
  boundary lattice points together) and whose bond factors carry the
  wheel-condition multipliers tracking the interior-point positions
  (higher-rank shuffle route).
\end{enumerate}
Direct multi-chart assembly with only the chart-wise $Y^+(\widehat{\fgl}_1)$
data misses the interior-point contributions to the Hall product.
\end{proposition}

\begin{proof}[Sketch]
The chart-wise single-vertex polyvectors generate only the
boundary-vertex portion of the full toric quiver shuffle algebra; the
interior-vertex portion is generated by compositions extending into the
$G$-McKay inertia (Bridgeland--King--Reid 2001 \texttt{arXiv:math/9908027}
Thm.~1.2).  The Rapcak--Soibelman--Yang--Zhao toric theorem (2020
\texttt{arXiv:2001.10549} Thm.~B) identifies the full quiver shuffle
data with $\CoHA_{\mathrm{DM}}(X_\Sigma)$ when the interior vertices are
included; the McKay/orbifold route recovers the interior through the
$[\C^3/G]$-orbifold chart and the BKR equivalence, and the higher-rank
shuffle route recovers it by enlarging the single-vertex chart shuffle
to the full $(Q_X, W_X)$-shuffle with interior-vertex weights inserted.
\end{proof}

% ============================================================================
\section{Local-to-global theorem for zero-interior-lattice-point diagrams}
\label{sec:gluing-p2-local-to-global}
% ============================================================================

The combination of \S\S\ref{sec:gluing-p2-conifold}--\ref{sec:gluing-p2-compact-cycle}
yields the closure theorem of Stratum~i: for a local toric CY$_3$ whose
toric diagram has no interior lattice points, the chart-wise \v{C}ech
assembly is exact and reproduces the Davison--Meinhardt critical CoHA.

\begin{theorem}[Local-to-global for zero-interior-point toric CY$_3$]
\label{thm:local-to-global-zero-interior}
Let $X_\Sigma$ be a smooth local toric Calabi--Yau threefold whose toric
diagram $\Delta_\Sigma$ has zero interior lattice points.  Then:
\begin{enumerate}[label=\emph{(\arabic*)}, leftmargin=2em]
\item The multi-chart \v{C}ech factorisation-algebra assembly
\[
  \cF_X \;=\; \mathrm{hocolim}_{\sigma \in \Sigma(3)}\,
    \cF_\sigma
    \;\big|\;
    \bigl\{T_{\sigma \sigma'} : \cF_\sigma \leftrightarrow \cF_{\sigma'}\bigr\},
\]
with $\cF_\sigma = \Phi^{FA}_3(\Gamma_\sigma)$ the Ginzburg dga of the
chart Jordan-triple $(Q_{\mathrm{Jor}}, W_\sigma)$ and $T_{\sigma \sigma'}$
the Fourier--Mukai transition with kernel $K_{\sigma \sigma'} = \cO_\Delta
\otimes \pi_1^*\cO_{\bP^1_\tau}(-1)$, is a well-defined $E_3$-factorisation
algebra on $X$.
\item The associated critical CoHA
\[
  \CoHA_{\mathrm{\check Cech}}(X)
  \;=\;
  \mathrm{hocolim}_{\sigma \in \Sigma(3)}\,
    Y^+(\widehat{\fgl}_1)_\sigma
    \;\big|\;
    \bigl\{K_{\sigma \sigma'}\bigr\}
\]
is isomorphic to the Davison--Meinhardt critical cohomological Hall
algebra
\(
  \CoHA_{\mathrm{DM}}(X) = \CoHA(Q_X, W_X)
\)
of the toric quiver with potential $(Q_X, W_X)$ attached to $\Sigma$.
\item The resulting algebra is the Drinfeld positive half of an affine
super-Yangian of the root system determined by $(Q_X, W_X)$, with rank
equal to the number of vertices of $Q_X$ and super structure controlled
by the parity of the base-orientation sign on each codimension-one face
of the fan.
\end{enumerate}
\end{theorem}

\begin{proof}[Sketch]
The \v{C}ech homotopy colimit is computed face by face: each
codimension-one face contributes an FM kernel~\eqref{eq:FM-kernel-conifold}
scaled by the multiplier of equation~\eqref{eq:FM-multiplier}; each
codimension-two face contributes a cocycle compatibility checked by
the Jacobian calculation of equation~\eqref{eq:jacobian-calculation}.
Zero interior lattice points means no codimension-three (interior)
corrections from the Neguț wheel condition, so each face contribution is
the direct Schiffmann--Vasserot shuffle kernel without higher-rank
adjustment.  The identification with $\CoHA_{\mathrm{DM}}(X)$ is the
Rapcak--Soibelman--Yang--Zhao toric theorem (2020
\texttt{arXiv:2001.10549} Thm.~A), which matches the chart-wise
Schiffmann--Vasserot data to the toric quiver BPS algebra.  The
super-structure on the output is the base-orientation cocycle of the
fan as computed in equation~\eqref{eq:jacobian-calculation}: for fans
with no internal edges (no compact $\bP^1$-bases), the cocycle is
trivial and the output is an even affine Yangian; for fans with one
internal edge (e.g.\ the conifold), the cocycle produces the super
structure $\widehat{\fgl}(1|1)$ of equation~\eqref{eq:conifold-master};
for fans with several internal edges, the cocycle produces the product
super structure of the corresponding super-Yangian.
\end{proof}

The class of local toric CY$_3$ covered by
Theorem~\ref{thm:local-to-global-zero-interior} includes:
\begin{itemize}
\item $\C^3$ (one chart, no overlaps, $\CoHA = Y^+(\widehat{\fgl}_1)$);
\item the resolved conifold (two charts, one overlap,
  $\CoHA = Y^+(\widehat{\fgl}(1|1))^{\mathrm{con}}$);
\item the banana threefold (three charts, three overlaps, three-loop fan);
\item the suspended pinch point (four charts, six overlaps);
\item generalised conifolds $X_{n, m}$ (two charts, one overlap with
  $\cO(-n) \oplus \cO(-m)$-twist).
\end{itemize}

The class \emph{not} covered by the theorem comprises exactly those local
toric CY$_3$ with one or more interior lattice points --- local $\bP^2$,
local $\bP^1 \times \bP^1$, local Hirzebruch $\bF_n$ for $n \geq 1$, local
del Pezzo $dP_k$ for $k \geq 1$.  These all carry compact $4$-cycles and
must pass through Proposition~\ref{prop:compact-cycle-obstruction}'s
McKay/orbifold or higher-rank shuffle resolutions.  Part III takes up
this resolution, replacing the direct chart-wise \v{C}ech assembly by the
BKR equivalence and the inertia-stack decomposition of the orbifold
route; the resulting picture extends to the full toric Stratum~i up to
$\Z_n$-orbifolds and recovers, in particular, the nine-arrow Beilinson
CoHA of local $\bP^2$ as the \v{C}ech assembly of the $[\C^3/\Z_3]$-chart
across its $\Z_3$-skew-group algebra.

\bigskip

The chart-wise assembly of Stratum~i is now closed.  Three invariants have
been extracted: the face-lattice cocycle with values in
$H^1(\Sigma, \underline{\mathrm{Pic}})$; the bond-factor presentation of
Galakhov--Li--Yamazaki on the toric diagram edges; the zero-interior-point
criterion that separates the direct-\v{C}ech-assembly diagrams from those
requiring McKay resolution.  Part III addresses the McKay route for
diagrams with one or more interior points, and Part IV takes up derived
Morita gluing for the Van den Bergh tilting presentation common to
both.  The master test case $K3 \times E$, which combines Stratum i with
Strata ii--iv, is resolved in Part VIII via the full four-stratum
decomposition.
